% --- INSTITUTIONAL TEMPLATE FED v7.4.9 ---
% Estándar de Reporteo APA 7 (Babel-Free)
\documentclass[12pt,spanish]{article}

\usepackage{graphicx}
\usepackage{geometry}
\usepackage{indentfirst} 
\usepackage{caption}
\usepackage{floatrow} 

% Forzar estilos de caption antes de cualquier carga
\floatsetup[figure]{capposition=top}
\floatsetup[table]{capposition=top}

\usepackage{booktabs}
\usepackage{float}
\usepackage{longtable}
\usepackage{threeparttable}
\usepackage{array}
\usepackage{calc}
\usepackage{amsmath, amssymb}
\usepackage{fancyhdr}
\usepackage{mathptmx} 
\usepackage[T1]{fontenc}
\usepackage{setspace} 
\usepackage{ragged2e} 
\usepackage{titlesec}
\usepackage{url}
\def\UrlBreaks{\do\/\do-\do\.\do\_\do\?\do\=\do\&\do\%}
\usepackage[hidelinks,breaklinks]{hyperref}

% --- Pandoc Compatibility Block (Iron Clad) ---
\providecommand{\tightlist}{%
  \setlength{\itemsep}{0pt}\setlength{\parskip}{0pt}}

\usepackage{xcolor}
\usepackage{fancyvrb}
\usepackage{framed}

% --- Code Highlighting (Pandoc Standard) ---
\AtBeginDocument{
  \definecolor{shadecolor}{RGB}{248,248,248}
}

% Compatibilidad Pandoc 3.x
\ifdefined\pandocbounded\else
  \newcommand{\pandocbounded}[1]{#1}
\fi

% --- Pandoc Citeproc (CSL) compatibility ---
\newlength{\cslhangindent}
\setlength{\cslhangindent}{1.27cm} 
\newlength{\csllabelwidth}
\setlength{\csllabelwidth}{3em}
\newcommand{\citeproctext}{} 
\newcommand{\citeprocitem}[2]{\item} 
\newenvironment{CSLReferences}[2] 
 {\begin{list}{}{
  \setlength{\leftmargin}{\cslhangindent}
  \setlength{\parsep}{0pt}
  \setlength{\itemsep}{6pt}
  \ifodd #1
   \setlength{\leftmargin}{\leftmargin + \csllabelwidth}
   \setlength{\itemindent}{-\csllabelwidth}
  \fi
  }
  \renewcommand{\bibitem}[2][]{\item} % Elimina los corchetes [ ]
 }
 {\end{list}}
\newcommand{\CSLBlock}[1]{#1\hfill\break}
\newcommand{\CSLLeftMargin}[1]{\parbox[t]{\csllabelwidth}{#1}}
\newcommand{\CSLRightInline}[1]{\parbox[t]{\linewidth - \csllabelwidth}{#1}\break}
\newcommand{\CSLIndent}[1]{\hspace{\cslhangindent}#1}

\makeatletter
\@ifundefined{paragraph}{\newcounter{paragraph}}{}
\@ifundefined{subparagraph}{\newcounter{subparagraph}}{}
\makeatother

% Solución al error 'No counter none defined' en longtable
\makeatletter
\newcounter{none}
\@ifpackageloaded{caption}{
  \DeclareCaptionType{none}
  \captionsetup[none]{name=Tabla}
}{}
\makeatother

\hypersetup{
  colorlinks=true,
  linkcolor=blue,
  filecolor=blue,
  citecolor=blue,
  urlcolor=blue
}

% --- Macros Institucionales ---
\newcommand{\fuente}[1]{
  \vspace{4pt}
  \begin{flushleft}
    \small
    \textit{Nota.} #1
  \end{flushleft}
}

% --- Localización Manual Determinista (Sin Babel) ---
\AtBeginDocument{
  \renewcommand{\tablename}{Tabla}
  \renewcommand{\figurename}{Figura}
  \renewcommand{\contentsname}{Índice}
  \renewcommand{\listtablename}{Índice de Tablas}
  \renewcommand{\listfigurename}{Índice de Figuras}
  \renewcommand{\abstractname}{Resumen}
  \renewcommand{\refname}{Referencias}
  
  \RaggedRight
  \setlength{\parindent}{1.27cm}
  \setlength{\parskip}{0pt}
}

\DeclareCaptionLabelSeparator{newline}{\\ \vspace{2pt}}
\captionsetup{
  position=top,
  justification=RaggedRight,
  singlelinecheck=false,
  labelfont=bf,
  textfont=it,
  labelsep=newline,
  font=small,
  skip=10pt
}
\captionsetup[figure]{name=Figura}
\captionsetup[table]{name=Tabla}

% Márgenes APA 7
\geometry{
  a4paper,
  margin=25.4mm,
  headheight=15mm
}

% --- Encabezados ---
\pagestyle{fancy}
\fancyhf{} 
\fancyhead[L]{\includegraphics[height=1.2cm]{/home/erick-fcs/Capital_Workstation/capital-workstation-libs/src/ecs_quantitative/reporting/assets/logos/logo_unl.png}}
\fancyhead[R]{\includegraphics[height=1.2cm]{/home/erick-fcs/Capital_Workstation/capital-workstation-libs/src/ecs_quantitative/reporting/assets/logos/logo_economia.jpg}}
\renewcommand{\headrulewidth}{0pt}
\fancyfoot[R]{\thepage}

\fancypagestyle{plain}{
  \fancyhf{}
  \fancyhead[L]{\includegraphics[height=1.2cm]{/home/erick-fcs/Capital_Workstation/capital-workstation-libs/src/ecs_quantitative/reporting/assets/logos/logo_unl.png}}
  \fancyhead[R]{\includegraphics[height=1.2cm]{/home/erick-fcs/Capital_Workstation/capital-workstation-libs/src/ecs_quantitative/reporting/assets/logos/logo_economia.jpg}}
  \fancyfoot[R]{\thepage}
  \renewcommand{\headrulewidth}{0pt}
}

% --- Títulos ---
\titleformat{\section}{\normalfont\normalsize\bfseries\centering}{\thesection}{1em}{}
\titleformat{\subsection}{\normalfont\normalsize\bfseries}{\thesubsection}{1em}{}
\titleformat{\subsubsection}{\normalfont\normalsize\bfseries\itshape}{\thesubsubsection}{1em}{}
\titlespacing*{\section}{0pt}{12pt}{6pt}
\titlespacing*{\subsection}{0pt}{12pt}{6pt}
\titlespacing*{\subsubsection}{0pt}{12pt}{6pt}

\newenvironment{referenciasAPA}{
  \setlength{\parindent}{-1.27cm}
  \setlength{\leftskip}{1.27cm}
  \setlength{\parskip}{6pt}
}{\par}

\usepackage{titlesec}
\usepackage{enumitem}\makeatletter\@ifpackageloaded{caption}{}{\usepackage{caption}}\AtBeginDocument{%
\ifdefined\contentsname
  \renewcommand*\contentsname{Table of contents}
\else
  \newcommand\contentsname{Table of contents}
\fi
\ifdefined\listfigurename
  \renewcommand*\listfigurename{List of Figures}
\else
  \newcommand\listfigurename{List of Figures}
\fi
\ifdefined\listtablename
  \renewcommand*\listtablename{List of Tables}
\else
  \newcommand\listtablename{List of Tables}
\fi
\ifdefined\figurename
  \renewcommand*\figurename{Figure}
\else
  \newcommand\figurename{Figure}
\fi
\ifdefined\tablename
  \renewcommand*\tablename{Table}
\else
  \newcommand\tablename{Table}
\fi
}\@ifpackageloaded{float}{}{\usepackage{float}}
\floatstyle{ruled}
\@ifundefined{c@chapter}{\newfloat{codelisting}{h}{lop}}{\newfloat{codelisting}{h}{lop}[chapter]}
\floatname{codelisting}{Listing}\newcommand*\listoflistings{\listof{codelisting}{List of Listings}}\makeatother\makeatletter\makeatother\makeatletter\@ifpackageloaded{caption}{}{\usepackage{caption}}
\@ifpackageloaded{subcaption}{}{\usepackage{subcaption}}\makeatother

\begin{document}

\begin{center}
    {\bfseries \Large UNIVERSIDAD NACIONAL DE LOJA}\\[0.3cm]
    {\bfseries \large FACULTAD JURÍDICA, SOCIAL Y
ADMINISTRATIVA}\\[0.2cm]
    {\bfseries CARRERA DE ECONOMÍA}\\[1.5cm]
    
    {\bfseries \large POLÍTICA ECONÓMICA}\\[1cm]
    
    {\bfseries \Large Aprendizaje Autónomo}\\[1.5cm]
\end{center}

\noindent {\bfseries Estudiante:} Erick Fabricio Condoy Seraquive \\
\noindent {\bfseries Docente:} Econ. Maria Gabriela Moreno Hurtado \\
\noindent {\bfseries Fecha:} 20 de abril de 2026 \\

\vspace{1cm}

\doublespacing
\setlength{\parindent}{1.27cm}

\section{Introducción}\label{introducciuxf3n}

La distinción entre una política económica técnicamente fundamentada y
un conjunto disperso de decisiones reactivas no es meramente semántica;
radica en la aplicación rigurosa de la \textbf{teoría de la política
económica} como disciplina frente a la simple \textbf{praxis} política.
Mientras que las decisiones reactivas suelen ser respuestas
\emph{ad-hoc} a presiones externas, una política fundamentada se
estructura sobre una jerarquía lógica de fines, objetivos e
instrumentos, tal como lo define Cuadrado Roura (2014).

\section{Desarrollo}\label{desarrollo}

\subsection{Diferencia entre Economía y Política
Económica}\label{diferencia-entre-economuxeda-y-poluxedtica-econuxf3mica}

La \textbf{Economía} (como ciencia positiva) se encarga de explicar los
fenómenos tal cual son, buscando relaciones de causalidad. En contraste,
la \textbf{Política Económica} es la aplicación de dicha ciencia para
modificar la realidad. Según Tinbergen (1952), esto implica la
``variación deliberada de medios para alcanzar objetivos''.

En el caso del \textbf{Acuerdo SAF 2024 con el FMI}, observamos una
política fundamentada porque no se limita a ``reaccionar'' a la falta de
liquidez, sino que establece un marco de ordenación (dimensión positiva)
basado en diagnósticos de sostenibilidad fiscal, traduciéndolos en
medidas normativas (lo que ``debe ser'') para garantizar la estabilidad
monetaria.

\subsection{Dimensión Positiva y
Normativa}\label{dimensiuxf3n-positiva-y-normativa}

Una política reactiva se queda en lo normativo del momento
(``necesitamos dinero ahora''). Una política técnica equilibra: 1.
\textbf{Dimensión Positiva}: El diagnóstico del déficit y la brecha de
financiamiento. 2. \textbf{Dimensión Normativa}: La decisión de
priorizar la estabilidad de la dolarización frente a objetivos de gasto
inmediato.

\subsection{Jerarquía de Fines, Objetivos e
Instrumentos}\label{jerarquuxeda-de-fines-objetivos-e-instrumentos}

La política técnica se diferencia por su estructura: - \textbf{Fines}:
Sostenibilidad institucional y blindaje de la dolarización. -
\textbf{Objetivos}: Consolidación fiscal (reducción del déficit
primario) y acumulación de reservas. - \textbf{Instrumentos}: Incremento
del IVA al 15\%, focalización de subsidios y establecimiento de pisos de
gasto social.

Esta jerarquía permite clasificar las políticas adecuadamente. Por
ejemplo, el Acuerdo SAF es una política de \textbf{proceso} (estabilidad
coyuntural) y de \textbf{ordenación} (reformas estructurales), lo que le
otorga coherencia interna, a diferencia de las medidas del caso
``Andesur'', que mezclan instrumentos (crédito) con aspiraciones sin una
ruta clara.

\section{Conclusión}\label{conclusiuxf3n}

Lo que distingue a una política fundamentada es su capacidad de
\textbf{jerarquización y justificación teórica}. El Acuerdo SAF 2024, a
pesar de sus costos sociales, responde a un diseño técnico que busca la
estabilidad de largo plazo mediante la manipulación de instrumentos
cuantitativos y cualitativos precisos. Una decisión reactiva, en cambio,
carece de esta trazabilidad y suele generar conflictos no previstos
entre objetivos al no considerar las restricciones presupuestarias y de
mercado.

\begin{center}\rule{0.5\linewidth}{0.5pt}\end{center}

\fuente{Cuadrado Roura (2014), FMI (2024) y análisis institucional del Nodo PE.}

\end{document}
